% Class-compatibility preamble retained for a reader excerpt of
%   Dave Witte Morris, "A Cohomological Proof that Real Representations of
%   Semisimple Lie Algebras Have Q-Forms", SIGMA 11 (2015) 034,
%   DOI 10.3842/SIGMA.2015.034, CC BY 4.0.
%
% The source pins the publisher class sigma, which is not part of the delivery
% baseline, so this file supplies the amsart baseline and then copies the author's
% own declaration block mechanically and unchanged:
%   - pinned-original-source.tex lines 11-31: the author's DeclareMathOperator
%     list, the equation numbering and the author's newtheorem declarations;
%   - pinned-original-source.tex lines 35-43, which the author writes immediately
%     after \begin{document} in the publisher layout: the author's own \Res,
%     \ResOne and \Lie macros and \allowdisplaybreaks.
% The publisher header commands of the sigma class (\FirstPageHeading,
% \ArticleName and the like) and the macro \PaperNumber they require are not
% reproduced, because this excerpt carries neither the journal front matter nor
% anything else that uses them.
% No mathematical macro, statement, symbol or assumption is introduced, changed
% or removed by this file.
\documentclass[10pt,a4paper]{amsart}
\usepackage{amssymb,amsmath,amsthm}
\usepackage{enumerate}
\usepackage{hyperref}

\DeclareMathOperator{\Cliff}{Clif\/f}
\DeclareMathOperator{\Gal}{Gal}
\DeclareMathOperator{\Inn}{Inn}
\DeclareMathOperator{\Mat}{Mat}
\DeclareMathOperator{\Out}{Out}
\numberwithin{equation}{section}

\newtheorem{prop}{Proposition}[section]
\newtheorem{cor}[prop]{Corollary}
\newtheorem{lem}[prop]{Lemma}
{\theoremstyle{definition}
\newtheorem{defn}[prop]{Definition}
\newtheorem{notation}[prop]{Notation}
\newtheorem{rem}[prop]{Remark}
\newtheorem{ques}[prop]{Question}
\newtheorem{eg}[prop]{Example}
\newtheorem{erratum}[prop]{Erratum}}

\newcommand{\Res}[1]{\mathop{\mathbf{R}_{#1}}}
\newcommand{\ResOne}[1]{\mathop{\mathbf{R}^{(1)}_{#1}}}
\newcommand{\Lie}[1]{\mathfrak{\lowercase{#1}}}
\allowdisplaybreaks
\usepackage{cleveref}
\crefname{originalremark}{original remark}{original remarks}
\Crefname{originalremark}{Original remark}{Original remarks}
\makeatletter
\newlabel{OnlyZStar}{{2.4}{source}{}{sourcepointer-1}{}}
\newlabel{OnlyZStar@cref}{{[originalremark][0][]2.4}{[1][0][]source}{}{}{}}
\makeatother
\crefname{originallemma}{original lemma}{original lemmas}
\Crefname{originallemma}{Original lemma}{Original lemmas}
\makeatletter
\newlabel{RunivIff}{{2.1}{source}{}{sourcepointer-2}{}}
\newlabel{RunivIff@cref}{{[originallemma][0][]2.1}{[1][0][]source}{}{}{}}
\makeatother
\crefname{originaldefinition}{original definition}{original definitions}
\Crefname{originaldefinition}{Original definition}{Original definitions}
\makeatletter
\newlabel{RunivIff-DivAlg}{{3}{source}{}{sourcepointer-3}{}}
\newlabel{RunivIff-DivAlg@cref}{{[originaldefinition][0][]3}{[1][0][]source}{}{}{}}
\makeatother
\title{Every absolutely simple Lie algebra of type E8, F4 or G2 over Q is R-universal}
\author{Dave Witte Morris}
\date{doi:10.3842/SIGMA.2015.034; CC BY 4.0}
\begin{document}
\maketitle
\noindent\textit{Editorial note.} All delivered mathematical text is the author's own, copied from the pinned publisher source sources/algebra-fl02-a-cohomological-proof-that-real-re/pinned-original-source.tex; no statement or proof step was rewritten, reordered or removed.\par\medskip
\noindent\textit{Editorial note.} Preamble provenance: the source pins the publisher class sigma, which is not part of the delivery baseline, so this card uses sources/algebra-fl02-a-cohomological-proof-that-real-re/retained-preamble.tex, which supplies the amsart baseline, copies the author's DeclareMathOperator list, equation numbering and newtheorem declarations from pinned source lines 11-31, and copies the author's own \textbackslash{}Res, \textbackslash{}ResOne and \textbackslash{}Lie macros from pinned source lines 35-43, which the author writes immediately after \textbackslash{}begin\{document\}. No macro was invented and no mathematical symbol was changed.\par\medskip
\noindent\textit{Editorial note.} Bibliography provenance: the counted statement and its proof contain no citation, so this card carries no bibliography. The source has a thebibliography environment at pinned source lines 954-1016 whose three entries are used elsewhere in the paper; nothing from it is needed here and no entry was reconstructed.\par\medskip
\noindent\textit{Editorial note.} External references are navigation pointers only; every deliberately kept step of the author's argument stays inside the counted spans. Their printed numbers are the author's own numbering under the author's newtheorem declarations, computed mechanically by the lane's scanner.\par\medskip
\noindent\textit{Allowed standard prerequisites (editorial).}\begin{itemize}
\item The classification of absolutely simple Lie algebras over Q by type, and the fact that types E8, F4 and G2 have no outer automorphisms and trivial centre. Used inside the counted proof at source lines 507-509; the classification frame is set up in the section the counted statement belongs to (source lines 378-400). Makes the star-action trivial on these types, which is the hypothesis of the closing step.
\item The criterion for R-universality through the Tits algebras being split. The author's remark labelled OnlyZStar at source lines 461-468 and the author's equivalence labelled RunivIff at source lines 415-436, both invoked inside the counted proof. Turns the triviality of the star-action into the conclusion that every Tits algebra is Q itself, hence the algebra is R-universal.
\end{itemize}\medskip
\ifdefined\hypertarget\hypertarget{sourcepointer-1}{}\fi
\noindent\textit{Source reference.} Original remark 2.4: sources/algebra-fl02-a-cohomological-proof-that-real-re/pinned-original-source.tex, source lines 462--462. Navigation pointer only: the counted proof invokes this remark in its second sentence; its content is declared as a named standard prerequisite and is not reproduced.\par\medskip
\ifdefined\hypertarget\hypertarget{sourcepointer-2}{}\fi
\noindent\textit{Source reference.} Original lemma 2.1: sources/algebra-fl02-a-cohomological-proof-that-real-re/pinned-original-source.tex, source lines 416--416. Navigation pointer only: the counted proof invokes the author's equivalence in its last sentence; the equivalence has its own card in this lane.\par\medskip
\ifdefined\hypertarget\hypertarget{sourcepointer-3}{}\fi
\noindent\textit{Source reference.} Original definition 3: sources/algebra-fl02-a-cohomological-proof-that-real-re/pinned-original-source.tex, source lines 428--428. Navigation pointer only: it is the third labelled item of that equivalence, cited by the counted proof once.\par\medskip
\begin{lem}
\label{EFGAllUniv}
Every absolutely simple Lie algebra of type $E_8$, $F_4$, or~$G_2$ over~$\mathbb{Q}$ is $\mathbb{R}$-universal.
\end{lem}

\begin{proof}
These types have no outer automorphisms, so the $*$-action must be trivial.
Furthermore, simple groups of these types have trivial center.
Therefore, it is immediate from Remark~\ref{OnlyZStar} that $D_{\Lie G, \mathbb{Q}}(\lambda) = \mathbb{Q}$ for every dominant
weight~$\lambda$.
So Lemma~\ref{RunivIff}(\ref{RunivIff-DivAlg}) implies that $\Lie G$ is $\mathbb{R}$-universal.
\end{proof}

% Bibliography excerpt retained for reader excerpts of
%   Dave Witte Morris, "A Cohomological Proof that Real Representations of
%   Semisimple Lie Algebras Have Q-Forms", SIGMA 11 (2015) 034,
%   DOI 10.3842/SIGMA.2015.034, CC BY 4.0.
%
% Every entry below is copied character for character from the author's own
% thebibliography environment of the pinned publisher source
% (pinned-original-source.tex lines 954-1016).  No entry was rewritten,
% reordered, merged, split or invented, and no citation key was changed.
\begin{thebibliography}{99}

\bibitem{Morris-RRepsQForms}
Morris D., Real representations of semisimple {L}ie algebras have {${\mathbb
  Q}$}-forms, in Algebraic Groups and Arithmetic, Tata Institute of Fundamental
  Research, Mumbai, 2004, 469--490.


\bibitem{PrasadRapinchuk-Prescribed}
Prasad G., Rapinchuk A., On the existence of isotropic forms of semi-simple
  algebraic groups over number f\/ields with prescribed local behavior,
  \href{http://dx.doi.org/10.1016/j.aim.2006.01.001}{\textit{Adv. Math.}} \textbf{207} (2006), 646--660.


\bibitem{Tits-RepLinIrred}
Tits J., Repr\'esentations lin\'eaires irr\'eductibles d'un groupe r\'eductif
  sur un corps quelconque, \href{http://dx.doi.org/10.1515/crll.1971.247.196}{\textit{J.~Reine Angew. Math.}} \textbf{247} (1971),
  196--220.



\end{thebibliography}
\end{document}
